\section{The Responsibility Gap, Which No Experiment Closes}
\label{sec:11.9}
An AI system's ethical behavior fails, and the parties who might answer for it — the team that designed it, the team that trained it, the organization that deployed it, the person who used it — each held only partial control over what it did. Traditional responsibility ascription assumes an operator who could have known what a system would do and could have done otherwise. A system built from training data nobody fully curated, tuned through a process nobody fully specified, and deployed into situations nobody fully anticipated breaks both conditions at once, for every party in the chain simultaneously. The philosopher Andreas Matthias named this the responsibility gap in 2004: a genuinely learning system creates a case existing law and moral theory were not built to handle \autocite{matthias2004responsibility}.

Filippo Santoni de Sio and Giulio Mecacci's 2021 refinement found not one gap but four — in culpability, in moral accountability, in public accountability, and in what they call active responsibility, the duty to prevent harm before it happens rather than answer for it after \autocite{santonidesio2021four}. The causes differ: technical unpredictability produces one, organizational diffusion another, and a mismatch between what designers expect and what the public relying on the system expects a third. Closing one does not close the others; a fix that makes a decision more explainable does nothing about the dozen teams that touched the training pipeline before it was made.

Legislation has answered part of this and left the part this book is about. The European Union's recast product liability directive, adopted in 2024, brings software expressly within the definition of a product and applies strict liability to defects in it; member states must transpose it by 9 December 2026, and its rules reach products placed on the market after that date \autocite{eu2024productliability}. Two features decide how far it travels. It carries the instrument built for opacity — a duty to disclose evidence, and rebuttable presumptions of defectiveness and of the causal link where technical complexity makes a claimant's proof excessively difficult — so the mechanism designed for systems nobody outside can inspect is now law rather than a proposal. And it recovers a closed list of damage: death, personal injury, damage to property, and the destruction or corruption of data. Pure economic loss is outside it, and so is nearly every harm this book has been about — the discriminatory score, the refusal nobody could appeal, the person selected for a strike.

The proposal that would have covered those was withdrawn. The Commission's 2022 AI Liability Directive would have extended a fault-based route, carrying the same presumption of causality, beyond defect and beyond the product regime's list of damage. Member states and industry could not agree on its terms; the Commission withdrew it in February 2025 and struck it from the legislative docket that October \autocite{europeancommission2022liability}. So the disclosure duty and the presumption hold where the harm is physical and the defendant is an economic operator in a chain built for goods, and lapse where the harm is to a right, the loss is economic, or the operator is a state. What is left in that second region is the gap Matthias described: liability distributed across designer, trainer, deployer, and user, governed by whatever each national law happens to say about a case none of them were written to anticipate. And a compensation regime of either kind answers culpability alone. Moral accountability, public accountability, and the duty to prevent harm before it happens are where the paragraph above left them.

A related, narrower case is what the law does once a system's own legal status is contested, sentient enough that the harm it causes or suffers might not reduce cleanly to someone else's liability. This problem is prior to that one and separate from it — it holds even for a system nobody claims is sentient, wherever control over its behavior is split across more than one party. Responsibility should track actual causal contribution. What nobody has produced is a workable way to measure that contribution after the fact, from outside a system whose own designers may not fully know why it did what it did.

\runin{The near-term work}
There is no experiment here. What would make it a research problem is a method for apportioning causal contribution across a training pipeline after a harm; absent one, the answer is allocated rather than measured, which is a drafting question. Two drafts exist. One was enacted and confined to defect; the other was negotiated, withdrawn, and left on the record with its terms intact. Neither was rejected for being unworkable, and the second's are available to whoever drafts the next one.

It belongs at the end because every other entry in this chapter asks for the same two things, and this is where one of them exists. Access to evidence a party would rather withhold, backed by an adverse presumption when it is withheld, is an enforceable duty now — for defect claims, against economic operators, over physical harm. That is the reach, and the confinement is the finding: outside it, a party outside the operator can be granted access, and publish what it finds, and still have produced nothing anybody must act on until somebody is answerable for it.
